% Part: lambda-calculus
% Chapter: church-rosser
% Section: beta-reduction

\documentclass[../../../include/open-logic-section]{subfiles}

\begin{document}

\olfileid[pt]{lam}{cr}{b}

\olsection{$\beta$-redução}

\begin{lem}\ollabel{lem:one-par}
  Se $M \bredone M'$, então $M \bredpar M'$.
\end{lem}
\begin{proof} Se $M \bredone M'$, então $M$ é
    $(\lambd[x][N])Q$, $M'$ é $\Subst{N}{Q}{x}$, para alguns
    $x$, $N$ e~$Q$. Como $N \bredpar N$ e $Q \bredpar Q$ por
    \olref[pb]{thm:refl}, temos imediatamente $(\lambd[x][N])Q
    \bredpar \Subst{N}{Q}{x}$ por
    \olref[pb]{defn:bredpar}\olref[pb]{defn:bredpar4}.
\end{proof}

\begin{lem}\ollabel{lem:par-red}
  Se $M \bredpar M'$, então $M \bred M'$.
\end{lem}

\begin{proof} Por indução na !!{derivation} de $M \bredpar M'$.
  \begin{enumerate}
  \item A última regra é \olref[pb]{defn:bredpar1}: então $M$ e $M'$ são
    simplesmente $x$, e $x \bred x$.
  \item A última regra é \olref[pb]{defn:bredpar2}: $M$ é
    $\lambd[x][N]$ e $M'$ é $\lambd[x][N']$ para alguns $x$, $N$, $N'$,
    onde $N \bredpar N'$. Pela hipótese de indução, temos $N \bred N'$.
    Então $\lambd[x][N] \bred \lambd[x][N']$ (pela mesma série de contrações
    $\bredone$ que em $N \bred N'$).
  \item A última regra é \olref[pb]{defn:bredpar3}: $M$ é $PQ$ e $M'$ é
    $P'Q'$ para alguns $P$, $Q$, $P'$, $Q'$, onde $P \bredpar P'$ e
    $Q \bredpar Q'$. Pela hipótese de indução, temos $P \bred P'$ e
    $Q \bred Q'$. Logo, $PQ \bred P'Q'$ pela sequência de redução
    $P \bred P'$ seguida pela redução $Q \bred Q'$.
  \item A última regra é \olref[pb]{defn:bredpar4}: $M$ é
    $(\lambd[x][N])Q$ e $M'$ é $\Subst{N'}{Q'}{x}$ para alguns $x$,
    $N$, $M'$, $Q$, $Q'$, onde $N \bredpar N'$ e $Q \bredpar Q'$.
    Pela hipótese de indução, obtemos $Q \bred Q'$ e $N \bred N'$. Assim,
    $(\lambd[x][N])Q \bred \Subst{N'}{Q'}{x}$ por $N \bred N'$, seguido
    de $Q \bred Q'$ e, finalmente, da contração de $(\lambd[x][N'])Q'$ a
    $\Subst{N'}{Q'}{x}$.
  \end{enumerate}
\end{proof}

\begin{lem}\ollabel{lem:str}
  $\bred$ é a menor relação transitiva que contém $\bredpar$.
\end{lem}

\begin{proof}
  Seja $\xred$ a menor relação transitiva que contém $\bredpar$.

  $\bred \subseteq \xred$: suponha $M \bred M'$, isto é, $M \ident M_1
  \bredone \dots \bredone M_k \ident M'$. Por \olref{lem:one-par}, $M
  \ident M_1 \bredpar \dots \bredpar M_k \ident M'$. Como $\xred$
  contém $\bredpar$ e é transitiva, $M \xred M'$.

  $\xred \subseteq \bred$: suponha $M \xred M'$, isto é, $M
  \ident M_1 \bredpar \dots \bredpar M_k \ident M'$. Por
  \olref{lem:par-red}, $M \ident M_1 \bred \dots \bred M_k \ident M'$.
  Como $\bred$ é transitiva, $M \bred M'$.
\end{proof}

\begin{thm}\ollabel{thm:cr}
  $\bred$ satisfaz a propriedade de Church--Rosser.
\end{thm}

\begin{proof}
  Imediato de \olref[dap]{thm:str}, \olref[pb]{thm:cr} e \olref{lem:str}.
\end{proof}

\end{document}

